% !TeX root = ../../../main.tex
\subsection{Clinical Reliability, Validation, and Error Analysis}
\label{ch06:sec:additional_validation}

\ChSixHeadingOnly

\subsubsection{Inter-annotator Agreement}
\label{ch06:sec:iaa}
To contextualize residual model--clinician discrepancies, we quantified agreement between two independent clinician annotators on five sampled case reports from T2S2, annotated under the same guidelines.
We report (i) model--annotator alignment for two leading LLM annotators against each clinician and (ii) clinician--clinician agreement using the same metrics as the main evaluation.


% !TeX root = ../../../main.tex
\begin{table}[htbp]
\caption{Inter-annotator agreement across reports.}\ChSixFloatLabel{ch06:tab:iaa_manual}
\ChSixTableBodyStart
\begin{ChSixFloatBody}

\centering
\small


\begin{adjustbox}{max width=\columnwidth}
\begin{tabular}{lccc}
\toprule
Annotator pair & Recall & Concordance & AULTC \\
\midrule
A1 vs.\ A2 (Manual) & 0.681 & 0.912 & 0.811 \\
\bottomrule
\end{tabular}
\end{adjustbox}

\end{ChSixFloatBody}
\end{table}







Tables~\ref{ch06:tab:iaa_manual} and~\ref{ch06:tab:iaa_models} contextualize the main results by quantifying annotator variability under the same guidelines and metrics.
Clinician--clinician agreement is imperfect even on the same case reports (Table~\ref{ch06:tab:iaa_manual}), reflecting inherent ambiguity in mapping narrative descriptions to temporally localized findings and timestamps.
Against this reference, the strongest LLM annotators achieve model--clinician agreement that is comparable to clinician--clinician agreement across the same metrics (Table~\ref{ch06:tab:iaa_models}), supporting that the leading performance levels in Table~\ref{ch06:tab:results} are meaningful relative to human variability rather than artifacts of the evaluation protocol.

\paragraph{Inter-annotator agreement on a GLP1-RA subset}
\label{ch05:glp_iaa}

\ChFiveRebuttal{To contextualize agreement with the single-expert reference set, we also analyzed a separate 140-case report subset of PMOA-TTS involving patients who took GLP-1 RA (Glucagon-like Peptide Receptor Agonists) medications. Two clinically trained annotators independently constructed event-time timelines for this subset using the same annotation framework. Using Annotator 1 as the reference and the same matching and temporal metrics used for the model evaluation, Annotator 2 achieved an event match recall of 0.811, temporal concordance of 0.798, and AULTC of 0.702.}

% !TeX root = ../../../main.tex

\subsubsection{Clinician Annotation Quality Review}

A clinician qualitative review of annotation quality on \texttt{sepsis-10} is described in Appendix~\ref{ch06:apd:clinician_quality_review}.

The reviewed annotation bundle included O1-preview as the OpenAI annotation; replacing GPT-5 which was evaluated separately in the main quantitative results in Table~\ref{ch06:tab:results}.

Inspecting the clinical quality of the annotations, we report the mean rank, top-1 fraction, top-1-LLM fraction, and quality level in Table~\ref{ch06:tab:sep10review}.
There is a clear performance ordering in the review, with manual annotations ranked highest, O1-preview a close second, followed by DeepSeek-R1 IQ1 and Llama 3.3.
While the Llama 3.3 annotations were rated acceptable a high percentage of the time (90\%), the other Llama and Mixtral approaches fail to achieve Acceptable performance in excess of 50\%.
Comparing DeepSeek-R1 IQ1 and Llama 3.3, the former has more variable clinician-rated quality (8 of 10 cases rated Acceptable or better), but maintains both higher event match recall and AULTC in the main quantitative evaluation (Table~\ref{ch06:tab:results} and Figure~\ref{ch06:fig:time-disc}).
The clinical reviewer also noted that some Llama 3.3 discrepancies for large time horizons, \emph{i.e.}, years, were large, which is corroborated by the long-horizon subpanels in Figure~\ref{ch06:fig:time-disc}.

% !TeX root = ../../../main.tex
\begin{table}[htbp]
\caption{Clinician annotation quality review on \texttt{sepsis-10}.}\ChSixFloatLabel{ch06:tab:sep10review}
\ChSixTableBodyStart
\begin{ChSixFloatBody}

\centering
\small


\begin{tabular}{lcccccc}
\toprule
Annotator & Mean Rank & Top-1 & Top-1-LLM & Excellent & $\geq$ Good & $\geq$ Acceptable \\
\midrule
Manual & 1.2 & 0.9 & -- & 1 & 1 & 1 \\
O1-preview & 1.5 & 0.6 & 0.9 & 0.9 & 1 & 1 \\
DeepSeek R1 IQ1 & 3.5 & 0.1 & 0.2 & 0.4 & 0.5 & 0.8 \\
L3.3 70B & 3.8 & 0 & 0 & 0 & 0.2 & 0.9 \\
L3.1 70B & 4.8 & 0 & 0 & 0 & 0.2 & 0.5 \\
Mixtral 8x7B & 5.4 & 0 & 0 & 0 & 0.1 & 0.2 \\
L3.1 8B & 6.5 & 0 & 0 & 0 & 0 & 0 \\
\bottomrule
\end{tabular}

\end{ChSixFloatBody}
\end{table}


\subsubsection{Sepsis Onset Timestamp Inference}
\label{ch06:sec:onset-inference}
In addition to extracting event--time pairs, we infer an explicit sepsis onset timestamp from narrative cues.
Given a case report, the onset procedure estimates the earliest time consistent with Sepsis-3 criteria relative to admission ($t=0$), producing two onset estimates based on qSOFA and SOFA-style reasoning when sufficient evidence is available.
We used GLM-4.6 (2-bit) for this inference.
The full prompt is provided in Appendix~\ref{ch06:apd:onset_prompt}.

\label{ch06:sec:onset_results}

On \texttt{sepsis-10}, SOFA-based onset matched clinician onset exactly in 6 of 9 valid cases where onset could be assessed, excluding one neonatal outlier for which neonatal SOFA would be more appropriate.
The median absolute error was 0 hours (mean $\approx 3.4$ hours), with a single moderate deviation of 24 hours.
In contrast, qSOFA-based onset tended to lag substantially (mean absolute error $\approx 355$ hours including the outlier and $\approx 126$ hours excluding it), consistent with retrospective reporting of qSOFA-related cues in case narratives rather than true physiologic delay, whereas SOFA-based onset closely aligned with clinician-annotated sepsis onset.

\subsubsection{Grounding and Consistency}
\label{ch06:sec:additional_validation_summary}

We further evaluated factual grounding, generation stability, and full-corpus consistency using analyses summarized here and detailed in Appendix~\ref{ch06:apd:validation_details}.
In the \texttt{sepsis-10} hallucination audit, 1,517 total extracted events from combined DeepSeek-R1 and Llama-3.3-70B-Instruct were reviewed through automated lexical grounding followed by manual adjudication of unresolved cases; no unsupported hallucinated events were identified (Table~\ref{ch06:tab:hallucination_audit} in Appendix~\ref{ch06:apd:hallucination_audit_results}).
On the same \texttt{sepsis-10} reports, multi-pass self-consistency across five random seeds showed stable extraction behavior across DeepSeek-R1, Llama 3.3, O1-preview, and GPT-5 (Table~\ref{ch06:tab:self_consistency} in Appendix~\ref{ch06:apd:self_consistency_results}).
Across the full 2,139-report corpus, Llama-3.3-70B-Instruct achieved strong agreement with DeepSeek-R1, with event match recall 0.79, concordance 0.87, and AULTC 0.83 (Table~\ref{ch06:tab:full_corpus_agreement} in Appendix~\ref{ch06:apd:full_corpus_agreement_results}).

\subsubsection{Annotation Error Analysis}
We qualitatively compared temporal annotations produced by several LLMs on a representative case report to illustrate characteristic differences in event extraction and timestamping (Supplementary Material, Section~\ref{ch05:apd:inter-llm-comparison}).
While all models recovered the major clinical trajectory and preserved temporal polarity, they varied in
event granularity, paraphrasing, and long-range temporal normalization.
These differences complement the aggregate quantitative results and highlight the impact of model behavior on the structure of extracted timelines.
Common model-specific error patterns, such as granularity variation, duplicate events, and differences in temporal normalization, are summarized in the Supplementary Material, Section~\ref{ch05:apd:inter-llm-comparison}, Table~\ref{ch05:tab:error_taxonomy}.
